\documentclass[11pt,a4paper]{article}
\usepackage{amsmath,amssymb}
\usepackage{geometry}
\usepackage{booktabs}
\usepackage{microtype}
\usepackage{parskip}
\geometry{margin=2.5cm}
\setlength{\parskip}{5pt}
\setlength{\parindent}{0pt}

\title{Symmetry Reduction in the Detailed Balance Checker:\\
\large Feasibility, Expected Speedup, and Correctness Conditions}
\author{}
\date{}

\begin{document}
\maketitle

% ============================================================
\section{What the method proposes}
% ============================================================

The detailed balance (DB) checker must verify
$T(s_i\!\to\!s_j)\,\pi(s_i) = T(s_j\!\to\!s_i)\,\pi(s_j)$
for every pair of states $(s_i,s_j)$ found by BFS.
For a component of $N$ states this is $N(N{-}1)/2$ pairs.
At \texttt{NGrid\,=\,3} ($N=504$) that is 126,756 pairs; at
\texttt{NGrid\,=\,6} ($N\approx100{,}000$) it is $\sim5\times10^9$,
which is infeasible.

The proposed speedup exploits the fact that for a translation-invariant
algorithm on a periodic lattice, all nine translated copies of a
configuration pair produce \emph{identical physics}.  If the pair
$(s,s')$ satisfies DB, then every translated copy $(g{\cdot}s,\,g{\cdot}s')$
automatically satisfies DB too, by a three-line argument:
\begin{align*}
  T(g{\cdot}s \to g{\cdot}s')\,\pi(g{\cdot}s)
  &= T(s\to s')\,\pi(s)
   = T(s'\to s)\,\pi(s')
   = T(g{\cdot}s'\to g{\cdot}s)\,\pi(g{\cdot}s').
\end{align*}
The outer equalities use the symmetry of $T$ and $\pi$; the middle uses
DB for the original pair.  The reverse direction holds identically:
a violation for $(s,s')$ implies a violation for every translated copy.
Therefore it suffices to check \emph{one pair per family of equivalent
pairs}.

For a $3\times3$ lattice with 9 translations ($|G|=9$), the 126,756
pairs reduce to $126{,}756/9 = 14{,}084$.  Adding the 8 rotations and
reflections of the square (D4, giving $|G|=72$) reduces this further to
${\approx}1{,}760$ pairs.

% ============================================================
\section{Where the runtime actually comes from}
% ============================================================

Understanding whether the speedup is real requires understanding the
current bottleneck precisely.

The checker pipeline is: (1) BFS to find all $N$ states; (2) for each
pair with non-zero $T$, build the symbolic DB expression; (3) deduplicate
expressions by hash and run \texttt{FastChecker} on each unique one;
(4) run \texttt{FullSimplify} on expressions \texttt{FastChecker} could
not resolve; (5) scan for violations.

Step (3) uses Mathematica's \texttt{Hash[]} on the symbolic expression
object.  This matches \emph{syntactically identical} expressions — same
atoms in the same positions — not merely mathematically equivalent ones.
G-equivalent pairs (e.g.\ the same cluster geometry translated to a
different lattice site) produce expressions that reference different site
indices, different neighbor lists, and different random-variable
assignments; they come out in different symbolic forms even though their
mathematical content is the same.  They therefore hash differently, and
the deduplication step does \emph{not} group them.

Consequence: the ${\sim}3{,}000$ unique expressions reported for
\texttt{NGrid\,=\,3} are genuinely distinct after all caching.  The
orbit reduction would reduce this count by the full factor of $|G|$, not
a fraction of it.

\textbf{Step (4), the FullSimplify calls, is the runtime bottleneck.}
For \texttt{vmmc\_lattice.wl} (polynomial conditions, no Erfc),
\texttt{FastChecker} resolves most expressions quickly, but a
fraction --- roughly 10\,\% of the unique expressions, or ${\sim}300$ at
\texttt{NGrid\,=\,3} --- require \texttt{FullSimplify}, each taking
${\sim}0.5$\,s.  At 300 calls $\times$ 0.5\,s, spread across 4 kernels,
total time is $\approx$\,3 minutes, consistent with observed runtimes.

% ============================================================
\section{Quantitative impact}
% ============================================================

All rows below refer to the 504-state connected component found at
\texttt{NGrid\,=\,3}.  The FullSimplify estimate assumes ${\sim}10\%$ of
unique expressions require it and ${\sim}0.5$\,s each.

\begin{center}
\renewcommand{\arraystretch}{1.25}
\begin{tabular}{lrrr}
\toprule
 & No symmetry & Translations ($|G|{=}9$) & D4+trans.\ ($|G|{=}72$) \\
\midrule
Pairs to check      & 126,756  & 14,084   & $\approx$1,760  \\
Unique sym.\ exprs  & $\sim$3,000 & $\sim$333 & $\sim$42    \\
FullSimplify calls  & $\sim$300   & $\sim$33  & $\sim$4     \\
Estimated runtime   & 3\,min   & $\sim$20\,s & $\sim$2\,s  \\
\bottomrule
\end{tabular}
\end{center}

\textbf{These estimates are approximate.}  The FullSimplify fraction
depends on the specific expressions generated, which varies with NGrid
and the seed configuration.  The table should be read as order-of-magnitude
guidance, not a performance guarantee.

For \texttt{NGrid\,=\,6} the reduction from $5\times10^9$ to
${\approx}1.7\times10^7$ pairs with $|G|{=}288$ is the difference between
infeasible and potentially tractable.  The absolute runtime at that scale
is harder to predict without profiling, since it depends on the fraction
of expressions needing FullSimplify, which may differ from the
NGrid\,=\,3 case.

% ============================================================
\section{Conditions for correctness, and their risks}
% ============================================================

The reduction is logically valid but depends on conditions that must be
verified, not assumed.  A failure in any one of them can cause the
checker to report PASS on a broken algorithm.

\textbf{1. The algorithm must actually be symmetric (critical).}
The entire reduction assumes $T(g{\cdot}s\to g{\cdot}s') = T(s\to s')$.
If this fails, families will contain pairs with different DB status, and
the checker will certify a violating pair based on a non-violating
representative.  \texttt{vmmc\_biased.wl} demonstrates this exactly: its
rightward bias means a pair involving a rightward cluster move violates
DB, but the $90^\circ$-rotated copy (an upward move, which is unbiased)
satisfies DB.  With D4 symmetry assumed, the upward pair becomes the
representative, and the checker incorrectly returns PASS.  Section 4
below describes how to verify symmetry from the source code.

\textbf{2. The connected component must be closed under the symmetry.}
Every symmetry image $g{\cdot}s$ of every state $s$ in the component
must also be in the same component.  If not, the effective symmetry group
for that component is smaller than $G$, and the claimed speedup factor is
wrong.  For \texttt{NGrid\,=\,3} with $|G|{=}9$, the component size
$504 = 56\times9$ is consistent with this, but it should be confirmed
programmatically, not inferred from arithmetic.

\textbf{3. BFS must not be short-circuited using canonical states.}
A tempting shortcut is to run BFS only from one representative state per
symmetry family and check only pairs of representatives.  This is wrong.
The canonical representative of $s$ is $g_0{\cdot}s$ and the canonical
representative of $s'$ is $g_1{\cdot}s'$ for (generically) different
$g_0\neq g_1$.  The transition $T(g_0{\cdot}s \to g_1{\cdot}s')$ is not
equal to $T(s\to s')$ in general.  BFS must run on all $N$ states, and
orbit reduction only controls which pairs are sent to the DB checker.

\textbf{4. Canonical representatives must be computed per pair.}
Some configurations are self-symmetric (a non-trivial $g$ fixes the
state $g{\cdot}s=s$), making their family smaller than $|G|$.  The
implementation must compute the canonical representative explicitly for
every pair, not assume all families have size $|G|$ and skip by that
factor.

\textbf{5. Rotation formulas must be correct for the lattice size.}
The $90^\circ$ rotation $(r,c)\mapsto(c,\,L{-}1{-}r)$ maps lattice sites
to lattice sites for any $L$, but an off-by-one error produces a mapping
that is not a true bijection.  This bug is silent: the code runs without
error but computes wrong orbit representatives, potentially skipping pairs
that were never actually certified.

% ============================================================
\section{Verifying position-independence from source code}
% ============================================================

The symmetry reduction can only be trusted if the algorithm file is
genuinely position-independent.  For \texttt{vmmc\_lattice.wl}-style
algorithms, this can be checked by reading three specific parts of the
source code.

\textbf{Seed selection.}  Look for the line that chooses the seed
particle.  It should be \texttt{RandomChoice[occupied]} where
\texttt{occupied} is the list of all non-empty sites, with no weighting
by position.  Any position-dependent weighting (e.g.\ preferring sites
near the lattice centre) would break translation invariance.

\textbf{Displacement selection.}  Look for the line that chooses the
proposed move direction.  For translation invariance it must be
\texttt{RandomChoice[\$displacements]} where \texttt{\$displacements}
is a fixed list that does not depend on the current state or on which
site was seeded.  For D4 invariance additionally, the list must be closed
under rotation and reflection: the uniform box
$\{(\pm1,0),(0,\pm1),(\pm1,\pm1)\}$ satisfies this; any list that
over-represents one direction (as in \texttt{vmmc\_biased.wl}) does not.

\textbf{Energy computation.}  Find the function that computes pair
interaction energy.  It should call \texttt{couplingJ[a,\,b,\,d2]} where
\texttt{d2} is the squared \emph{periodic} distance between the two
sites:
\[
  d^2 = \min(\Delta r,\;L{-}\Delta r)^2 + \min(\Delta c,\;L{-}\Delta c)^2,
\]
with $\Delta r = |r_1-r_2|$ and $\Delta c = |c_1-c_2|$.  The distance
formula must use the minimum-image convention (the \texttt{min} terms)
to be periodic-boundary-correct.  If the energy instead uses raw
coordinate differences $\Delta r = r_1-r_2$ (not wrapped), it will not
be translation-invariant at the lattice boundary.  For D4 invariance,
note that $90^\circ$ rotation swaps $\Delta r$ and $\Delta c$ and
negates one; since both appear squared and summed, $d^2$ is unchanged.

These three checks can be performed by reading fewer than 20 lines of
any algorithm file.  They are necessary but not sufficient: they confirm
that the \emph{stated logic} is symmetric, but not that the supporting
library functions (\texttt{\$neighborsD2}, \texttt{\$applyDir}, etc.)
implement that logic correctly.  Those library functions should be
independently verified once and their symmetry properties documented.

% ============================================================
\section{Is it worth implementing?}
% ============================================================

\textbf{At \texttt{NGrid\,=\,3} (current scale):} Yes, the speedup is
real.  The hash-based deduplication does not capture G-equivalence, so
the orbit reduction genuinely reduces FullSimplify calls from ${\sim}300$
to ${\sim}4$ with D4+translations, cutting 3 minutes to ${\sim}2$
seconds.  However, 3 minutes is already tolerable for a correctness
check.  The implementation requires satisfying five non-trivial
correctness conditions, and a bug in any one of them introduces silent
false negatives.  The benefit at this scale does not clearly outweigh the
implementation risk.

\textbf{At \texttt{NGrid\,=\,6} and beyond:} The checker is entirely
infeasible without orbit reduction.  $5\times10^9$ pairs cannot be
constructed, let alone evaluated, in any reasonable time.  The orbit
reduction to ${\sim}1.7\times10^7$ pairs is not an optimisation; it is
a prerequisite for the checker to run at all.  At this scale the
implementation is worth doing despite the complexity.

\textbf{Recommendation:} Do not implement for \texttt{NGrid\,=\,3}.
Implement --- with all five correctness conditions satisfied explicitly
--- when \texttt{NGrid\,=\,5} or higher is a target.

\end{document}
